\chapter{Background}


\section{Domain: Machine Learning Research Methodology}

The landscape of Machine Learning (ML) research and development is 
characterized by rapid iteration and experimentation. According to the 
AI Index Report 2023, AI-related publications have more than doubled since 
2015, reflecting the field's exponential growth~\cite{AIIndex2023}. Data 
scientists and ML engineers continuously explore different model architectures, 
algorithms, hyperparameters, and data preprocessing techniques to achieve 
optimal performance. This iterative process, while essential for innovation, 
introduces significant challenges concerning efficiency and reproducibility. 
A single experiment may require hours or days of computation, and researchers 
often repeat similar computations across multiple runs. The cumulative effect 
translates into substantial costs in GPU hours, cloud computing expenses, 
and researcher productivity~\cite{schubotz2022caching}.

To manage this complexity, standardized process models have emerged. 
CRISP-DM (Cross-Industry Standard Process for Data Mining), developed in 
1996, remains widely adopted~\cite{wirth2000crispdm}. However, its limitations 
for modern ML—particularly lack of emphasis on model monitoring and 
maintenance, led to CRISP-ML(Q), which extends CRISP-DM with quality assurance 
methodology tailored for ML applications~\cite{studer2021crispmlq}. This 
framework aligns with modern MLOps principles, which emphasize automation, 
continuous integration/deployment (CI/CD), and reproducibility throughout 
the ML lifecycle~\cite{kreuzberger2023mlops}.

Tools such as MLflow~\cite{zaharia2018mlflow} and MLXP~\cite{arbel2024mlxp} 
support experiment tracking by recording parameters, metrics, code versions, 
and artifacts for each run, enabling comparison and reproducibility. Despite 
these advances, reproducibility remains a pervasive challenge in ML 
research~\cite{semmelrock2023reproducibility}. Reproducibility ensures consistent 
results from an ML workflow under identical conditions, hinging on the 
``Holy Trinity'': input data, code and parameters, and the execution 
environment~\cite{schubotz2022caching}. Challenges arise from:

\begin{itemize}
    \item \textbf{Lack of Records and Hyperparameter Inconsistency:} 
          Unrecorded decisions or parameter changes hinder replication.
    \item \textbf{Data and Environment Volatility:} Dataset alterations, 
          framework updates, and inconsistent compute environments cause 
          divergent outcomes~\cite{idowu2024empirical}.
    \item \textbf{Inherent Randomness:} Random initializations, data 
          shuffling, and non-deterministic GPU operations make consistent 
          replication challenging~\cite{semmelrock2023reproducibility}.
\end{itemize}

These challenges are particularly acute during iterative experimentation, 
where researchers may run hundreds of configurations before identifying 
promising approaches. Without intelligent management of computational 
artifacts, substantial resources are wasted on redundant calculations. 
This observation motivates the central contribution of this thesis: a 
fingerprint-based caching framework that persists and reuses computational 
artifacts (extracted features and trained models) across iterative modeling 
phases while maintaining reproducibility through cryptographic configuration 
fingerprinting.

\section{Context: Computational Efficiency in Research}

Building on the iterative experimentation challenges described in 
Section 2.1, this section examines the computational efficiency bottlenecks 
that motivate intelligent caching frameworks.

The increasing complexity and scale of Machine Learning models have 
substantially raised computational demands. Training models, tuning 
hyperparameters, and running repeated experiments consume vast compute, 
memory, and time resources, creating bottlenecks in research productivity 
and driving up costs.

Several factors drive this resource consumption:

\begin{itemize}
\item \textbf{Large Datasets:} Massive datasets require intensive I/O and 
      memory usage.
      
\item \textbf{Complex Models:} Both deep neural networks and ensemble 
      methods (e.g., gradient boosting with thousands of trees, large 
      random forests) require substantial computational resources for 
      training and inference.
      
\item \textbf{Iterative Experimentation:} Frequent experiments often repeat 
      computations unnecessarily when modifying a subset of parameters while 
      keeping others constant.
      
\item \textbf{Hyperparameter Optimization:} Searching large hyperparameter 
      spaces is inherently compute-intensive.
      
\item \textbf{Rigorous Validation:} Cross-validation strategies (e.g., 
      k-fold, leave-one-out) multiply training costs by the number of folds.
\end{itemize}

While caching can help, traditional caching systems (e.g., general-purpose 
key-value stores like Redis or Memcached) often fail to meet ML-specific 
needs~\cite{krishna2025cache}. These systems are typically static and focus 
on reducing latency rather than adapting to experiment configuration changes 
or optimizing for ML-specific access patterns. Such inefficiencies can slow 
research progress and increase costs.

This thesis addresses these inefficiencies through a fingerprint-based 
caching framework specifically designed for iterative ML workflows. By 
encoding complete experiment configurations (including cross-validation 
fold identity) into cryptographic cache keys, the framework enables 
automatic result reuse while maintaining scientific validity and 
reproducibility.

\section{Technology: Configuration Fingerprinting and Caching}

\begin{center}
\colorbox{yellow!20}{%
\begin{minipage}{0.9\textwidth}
\textbf{Draft Note:} Some statistics and system-specific claims in this section 
are from papers I reviewed, and written down notes but haven't yet added to the bibliography. These 
need proper bib entries and source verification before final submission. 
Marked with \textcolor{red}{[TODO: Add bib entry]}.
\end{minipage}
}
\end{center}

Efficient ML experimentation depends on configuration management and intelligent 
result reuse. This section reviews technologies directly relevant to the 
fingerprint-based LOSO caching approach developed in this thesis.

\subsection{Configuration Fingerprinting}

Configuration fingerprinting uses cryptographic hash functions (MD5, SHA-256) 
to create unique identifiers for experiment configurations. By hashing a 
canonical representation of all parameters, these fingerprints enable 
deterministic identification of identical configurations and automatic 
cache invalidation when parameters change.

Modern MLOps tools employ this approach:
\begin{itemize}
\item \textbf{DVC} logs pipeline signatures combining dependency contents, 
      parameters, and commands~\cite{dvc2022inproceedings}
\item \textbf{MLflow} tracks experiment metadata and enables visual 
      comparison~\cite{zaharia2018mlflow}
\item \textbf{MLXP} uses Hydra-based configuration logging with Git 
      hashes~\cite{arbel2024mlxp}
\end{itemize}

\subsection{Hash-Based Caching in MLOps}

MLOps pipeline caching implements content-addressable storage: when pipeline 
inputs and code match previous runs (verified through hash comparison), 
cached outputs are reused~\cite{dvc2022inproceedings}. DVC's run cache, 
for example, combines dependency hashes with parameter configurations to 
create cache keys. When a match is found, execution is skipped and previous 
results are restored.

This approach provides reproducibility: identical configurations produce 
identical cache keys, ensuring cached results correspond to recorded 
parameters~\cite{schubotz2022caching}.

\subsection{Related Caching Approaches}

\begin{center}
\colorbox{red!20}{%
\begin{minipage}{0.9\textwidth}
\textcolor{red}{\textbf{TODO:}} Add bib entries for: vLLM/PagedAttention paper 
(60-80\% memory reduction claim), GATI learned caching paper (7.69$\times$ speedup), 
XRootD/XCache infrastructure reports (10.8 PB statistic). Papers are in my 
download folder but not yet in references.bib.
\end{minipage}
}
\end{center}

While this thesis focuses on hash-based caching for ML experiments, other 
caching paradigms exist in adjacent domains:

\textbf{Large Language Model (LLM) Inference Caching} addresses inference 
latency by caching attention key-value tensors during autoregressive generation. 
Systems like vLLM's PagedAttention reportedly reduce memory waste from 60--80\% 
to under 4\% and enable up to 24$\times$ higher throughput \textcolor{red}{[TODO: Add bib entry]}. 
Prefix caching can achieve up to 85\% latency reduction and 90\% cost savings 
for repeated system prompts \textcolor{red}{[TODO: Add bib entry]}. These techniques 
are specific to LLM inference and not directly applicable to experiment result caching.

\textbf{Learned Caching for Deep Networks} employs neural networks to predict 
hidden layer outputs, enabling layer skipping during inference. GATI, for 
example, achieves up to 7.69$\times$ latency reduction at 97\% accuracy 
\textcolor{red}{[TODO: Add bib entry]}. This represents a different caching 
philosophy—using ML to make cache decisions—but applies only to inference scenarios.

\textbf{High-Energy Physics (HEP) Data Caching} operates at petabyte scale 
for raw data files. Systems like XRootD/XCache achieve cache hit rates exceeding 
85\% by request count. The Southern California Cache infrastructure reportedly 
saved 10.8 PB of network traffic in a single year \textcolor{red}{[TODO: Add bib entry]}. 
However, these systems are optimized for data distribution, not experiment 
configuration management.

In contrast to these approaches, this thesis implements deterministic, 
hash-based caching specifically designed for cross-validation experiments, 
where cache validity depends on both configuration parameters and fold identity.

\subsection{Reproducibility Considerations}

Despite fingerprinting, perfect reproducibility remains challenging. Sources 
of non-determinism include random seed variation, software version differences, 
hardware-specific floating-point operations, and non-deterministic GPU 
algorithms~\cite{semmelrock2023reproducibility}.

\begin{center}
\colorbox{red!20}{%
\begin{minipage}{0.9\textwidth}
\textcolor{red}{\textbf{TODO:}} Verify 0.2--0.5\% variance statistic. Check if 
this is in Semmelrock 2023 or requires separate citation (possibly Pham et al. 
2020 on variance in deep learning systems).
\end{minipage}
}
\end{center}

Empirical studies suggest that even with controlled random seeds, small 
variance (approximately 0.2--0.5\%) may persist between runs due to hardware-level 
floating-point non-determinism \textcolor{red}{[TODO: Add bib entry]}. Configuration 
fingerprinting addresses controllable factors (seeds, software versions, data 
splits) while accepting that hardware-induced variance cannot be eliminated 
through software means. For caching purposes, cached results are considered 
valid if configuration fingerprints match, with the understanding that minor 
numerical differences may occur due to non-deterministic hardware operations.
\chapter{Related Works}
